% =============================================================================
% 2-claims template — standalone-compilable lifecycle stage contract
%
% Structure: Hypotheses → Matrix → Detail → Alignment
% Format: venue-FREE evidence inventory; no [primary] designation
%
% Claims is the pure claim/evidence ledger.
% Venue-specific framing (RQs, editor's chair) lives in pitch.
% =============================================================================
\documentclass[11pt]{article}
\usepackage[margin=1in]{geometry}
\usepackage{parskip}
\usepackage{booktabs}
\usepackage{tabularx}
\usepackage{hyperref}
\usepackage{xcolor}
\newcommand{\status}[1]{\textsc{#1}}
\newcommand{\supported}{\textcolor{green!50!black}{\status{supported}}}
\newcommand{\weak}{\textcolor{orange!80!black}{\status{weak}}}
\newcommand{\gap}{\textcolor{red!70!black}{\status{gap}}}
\newcommand{\pending}{\textcolor{blue!60!black}{\status{pending}}}
\newcommand{\needprobe}[1]{\textcolor{red}{\textbf{[NEED PROBE]} #1}}
\title{Claims: <Paper Title>}
\date{YYYY-MM-DD}
\begin{document}
\maketitle

% ─────────────────────────────────────────────────────────────────────────────
% 1. HYPOTHESES
%    Venue-neutral. RQs and editor's chair framing live in pitch.
% ─────────────────────────────────────────────────────────────────────────────
\section*{Hypotheses}

\begin{enumerate}
\item[\textbf{H1.}] <venue-neutral hypothesis statement>
\item[\textbf{H2.}] <hypothesis>
\item[\textbf{H3.}] <hypothesis>
\end{enumerate}

% ─────────────────────────────────────────────────────────────────────────────
% 2. CLAIM-EVIDENCE MATRIX
%    Index: one row per claim, status at a glance.
% ─────────────────────────────────────────────────────────────────────────────
\section*{Claim-Evidence Matrix}

% =========================================================
% Para [claims.ledger] Result -- claim / status at a glance
% =========================================================
\begin{small}
\begin{tabularx}{\textwidth}{@{}p{0.06\textwidth}X p{0.12\textwidth}@{}}
\toprule
ID & Claim & Status \\
\midrule
C1 & <claim, as the paper states it> & \supported \\
\addlinespace
C2 & <claim> & \weak \\
\addlinespace
C3 & <claim> & \gap \\
\bottomrule
\end{tabularx}
\end{small}

% ─────────────────────────────────────────────────────────────────────────────
% 3. PER-CLAIM DETAIL
%    One \subsection* per claim. Four slots per paragraph:
%    S1 = claim + verdict, S2 = verified statistic, S3 = interpretation, S4 = source + caveat
% ─────────────────────────────────────────────────────────────────────────────
\section*{Per-Claim Detail}

\subsection*{C1. <short title> (\supported)}

% =========================================================
% Para [claims.c1] Result -- <the point>
% =========================================================
%% ---- P1.S1 ----
<claim statement>.
%
%% ---- P1.S2 ----
<verified statistic, spec, N>.
%
%% ---- P1.S3 ----
<one-line interpretation>.
%
%% ---- P1.S4 ----
Source: <real file>; <caveat if any>.

% ... more claims: C2, C3, etc. ...
% For weak/GAP claims: state the gap and the route instead of a statistic.

% ─────────────────────────────────────────────────────────────────────────────
% 4. DISCUSSION-ONLY INTERPRETATIONS (optional)
%    Interpretive findings that belong in Discussion, not Results.
%    Explicitly marked to prevent creep into Results claims.
% ─────────────────────────────────────────────────────────────────────────────
\section*{Discussion-Only Interpretations}

\textbf{D1. <label> (interpretive):} <description>. No direct measurement; interpretive frame only.

% ─────────────────────────────────────────────────────────────────────────────
% 5. ROBUSTNESS (optional)
%    Reported as a design strength in Methods, not claimed as a finding.
% ─────────────────────────────────────────────────────────────────────────────
\section*{Robustness (reported in Methods, not claimed)}

<sensitivity analyses: clustering, alt specs, multiple testing, exclusion criteria>

% ─────────────────────────────────────────────────────────────────────────────
% 6. PENDING EVIDENCE
%    Probes/tasks not yet run. Dependency chain.
% ─────────────────────────────────────────────────────────────────────────────
\section*{Pending Evidence}

\textbf{<probe id> (<label>, <status>):} <what it will test>. <which claim it upgrades>.

\textbf{Exploratory (supplement):} <analyses held from main claims>.

% ─────────────────────────────────────────────────────────────────────────────
% 7. HYPOTHESIS-CLAIM ALIGNMENT (at the bottom — validates the whole ledger)
%    Maps hypotheses to claims. No venue framing.
% ─────────────────────────────────────────────────────────────────────────────
\section*{Hypothesis-Claim Alignment}

\begin{small}
\begin{tabularx}{\textwidth}{@{}p{0.05\textwidth} p{0.30\textwidth} p{0.13\textwidth} X@{}}
\toprule
H & Hypothesis & Claims & How claims test this hypothesis \\
\midrule
H1 & <hypothesis> & C1, C2 & <how the claims provide evidence for/against H1> \\
\addlinespace
H2 & <hypothesis> & C3 & <how> \\
\addlinespace
H3 & <hypothesis> & C4 & <how> \\
\bottomrule
\end{tabularx}
\end{small}

\end{document}
